I selected JLCPCB as the PCB manufacturer since they are capable of producing the PCBs within their capabilities, and I (PAST) have always used JLC to produce our PCBs.
\nl
Because of the BGA package size, I concluded that it was not feasible to solder those two components by hand. To place the BGAs, I consulted with Peter from AT\&MWA to discuss their capabilities and whether they can place the two components or not. I found that:
\begin{itemize}[noitemsep]
    \item They are capable of placing BGAs of up to 0.2mm diameter and 0.5mm pitch.
    \item I should aim for a solder mask dam of 4mil (0.102mm) so the board can be fabricated with solder mask around each pad. I can use 2mil if required.
    \item I should produce a few spare boards in case I want to stress one or two out during testing.
    \item The two ICs should be placed in around 30-40 minutes a board. This would costs around \$35-\$47AUD+GST per board.
    \item They will produce the stentil themselves, and they will recommend the best way to panelise the boards to suit their machines.
    \item I are able to place the rest of the components first with the BGAs being last.
\end{itemize}